\documentclass[11pt]{article}
\usepackage[margin=1in]{geometry}
\usepackage[T1]{fontenc}
\usepackage{lmodern,amsmath,booktabs,graphicx}
\usepackage[hidelinks]{hyperref}\usepackage{xurl}
\setlength{\parindent}{0pt}\setlength{\parskip}{6pt}\setlength{\emergencystretch}{3em}
\title{Predicted Target Pose as the Initial Value for LiDAR Rectangle Fitting}
\author{Pan Dynamics Collision Avoidance Research}\date{October 3, 2026}
\begin{document}\maketitle
\section{Implementation and measured result}
The new perception path obtains both relative center and body heading from the
NRMM state predicted to the current sample. That planar pose initializes a
local rectangle fit in three free variables, $(c_x,c_y,\psi)$. Current point
cloud geometry can correct all three variables; none is locked or penalized
against the observer prediction. Declared length and width remain fixed.
The fitted center then becomes the current target-position observation.

The full perception and estimation replay was executed on the same previously
recorded 15 m and 20 m following-vehicle datasets used in earlier studies.
Actual YOLO26x inference, native LiDAR association, native ego preprocessing
and the committed MATLAB NRMM runtime were used. Each recording contains
300 evaluation frames at 0.02 s. The parked vehicles were already hidden in
these captures; no CARLA server was contacted for this work.

\begin{center}\scriptsize
\begin{tabular}{@{}llrrrrr@{}}\toprule
 & & \multicolumn{3}{c}{After 2 seconds} & \multicolumn{2}{c}{All frames} \\
Data & Fit method & Raw pos. & Est. pos. & Est. vel. & Est. pos. & Est. vel. \\
 & & m & m & m/s & m & m/s \\\midrule
15 m & Previous motion heading & 0.0674 & 0.0946 & 0.3049 & 0.1314 & 0.8195 \\
15 m & Fixed NRMM heading & 0.0810 & 0.1274 & 0.4471 & 0.1557 & 0.8663 \\
15 m & Predicted pose initialization & 0.2452 & 0.2648 & 0.4430 & 0.2354 & 0.8478 \\
20 m & Previous motion heading & 0.0661 & 0.1430 & 0.4997 & 0.6948 & 2.6381 \\
20 m & Fixed NRMM heading & 0.0719 & 0.1018 & 0.3782 & 0.1476 & 0.8506 \\
20 m & Predicted pose initialization & 0.0639 & 0.0637 & 0.1619 & 0.1155 & 0.7775 \\
\bottomrule\end{tabular}\end{center}

The 20 m result improves relative to the fixed-heading feedback implementation,
including its 0.3 s camera interruption. The 15 m position result regresses.
The pose-initialization API therefore remains optional. The comparison evaluates
the complete new local fitter: freeing translation also changes the earlier
face-anchored placement problem, so the gain cannot be attributed exclusively
to adding position to the seed. Both recordings were already seen; these are
conditional measured findings, not independent validation.

\section{Prediction, fitting and update order}
Let the runtime's state predicted from measurements through $k-1$ be available
at $t_k$. Its relative target position is already in current ego body axes.
The initialization for the measured point cloud is
\[
 z_k^0=\begin{bmatrix}\hat\rho_{x,k}^{-}&\hat\rho_{y,k}^{-}&
 \hat\psi_{C/E,k}^{-}\end{bmatrix}^{\!T},\qquad
 \hat\psi_{C/E}^{-}=\operatorname{atan2}(\hat q_y^{-},\hat q_x^{-})-\hat\beta_C^{-}.
\]
The native CARLA adapter changes the signs of the second position coordinate
and relative heading to match CARLA planar coordinates. It uses
\path{targetEstimate.relativePosition} and \path{targetRelativeHeading}
directly, without an estimated-world-yaw conversion.

The order is prediction query, current-cloud fit, then the existing
\path{onlineNrmmTrackingRuntime} \texttt{sample} action. The prior state time
must equal the frame time and its latest input time must strictly precede that
frame. Audit comparisons verify that each supplied predicted position and
heading equal the runtime's current published pre-update state. No repeated
estimator correction or same-frame prediction/fit feedback cycle is performed.
Both geometric fitting stages precede the single measurement update.

A missing detection is passed as a missing observation. The runtime continues
its prediction, but its predicted center is never passed back as a synthetic
measurement. The normal and interrupted runs have identical first 150
measurements and estimator states, exactly 15 missing camera detections at
frames 150--164, and reacquisition at frame 165.

There is no previous target pose at first acquisition. For this fixture, native
GNSS ego velocity supplies the existing co-moving initial heading assumption.
The first center is initialized from current point-cloud face placement, after
which all three pose variables are refined. This is labeled initialization;
it is not a converged target-state prediction or an arbitrary-acquisition solution.

\section{Geometry objective and local numerical method}
For a point $p_i$, define the rectangle-local coordinate
$u_i=R(\psi)^T(p_i-c)$ and the signed axis overflow
$e_i=|u_i|-[L/2,W/2]^T$. The optimized geometry cost is
\[
 J(c,\psi)=\frac1N\sum_i\left[
 \min\{\min_j |e_{ij}|^2,0.25\}
 +20\|\max(e_i,0)\|^2\right]+J_{\rm sym}.
\]
The boundary cap is in square metres. The residual is the existing nearest-face
closeness plus containment construction, rather than a likelihood for a
statistically exact vehicle-surface model. The dimensions are 5.0 by 1.9 m.

Symmetric placement uses the existing approximately symmetric visibility
assumption. At the beginning of each fitting stage, point projections in the
seed orientation select the most covered axis, if its support fraction is at
least 0.8. For that axis, $J_{\rm sym}$ is the squared midpoint of the two
support quantiles of $u_i$; otherwise it is zero. The support axis stays fixed
within that local solve. Thus the seed also initializes support selection; no
position or heading regularization term is added. The first stage uses 2 and
98 percent support quantiles, removes points outside the rectangle by more
than 0.15 m, and the refinement uses the first fitted pose and zero-percent
support quantiles on retained points.

A small NumPy damped Gauss--Newton implementation minimizes these residuals
using central-difference derivatives. Angular steps are scaled by rectangle
corner displacement. A 0.5 m scaled limit bounds each numerical step, not the
total correction; there is no hard yaw window around the prediction.
The iteration limit is 60 per stage, and convergence diagnostics are retained.
A finite, improved result at the limit is not a claim of global optimality.
No new SciPy or other numerical dependency is introduced.

Tests with two visible faces recover a known rectangle from biased centers and
headings. A separate partial-face test retains different tangential positions
from different seeds, correctly exposing missing geometric information.
Initialization can select a local solution, but cannot make an unobserved
direction observable.

\section{Initialization stress and the 15 m regression}
Cached actual detections and native clouds were used for two initialization
stress cases without changing the estimator equations or gains. With a
20-degree initial target heading error, the 20 m pose-initialized fit reaches
0.0661 m position and 0.1725 m/s velocity RMSE after 2 s. The preceding locked
heading experiment gave 0.6761 m and 2.5280 m/s. With a 90-degree error the new
post-2-second result is 0.1140 m and 0.7355 m/s, but the full-record position
RMSE remains 1.3167 m. Improved later recovery does not remove the large
acquisition error or establish global recovery.

The 15 m result demonstrates an objective-model limitation. At prototype frame
200, the first-stage objective evaluated at actor-truth pose is 0.12359, while
the fitted pose has lower cost 0.09782 despite about 0.47 m lateral displacement
and a 6.17-degree heading. The initial lateral support fraction is 0.786, below
the pre-existing 0.8 symmetry threshold. Refinement reduces the displacement
but does not eliminate the bias. Five sampled frames reproduce the same
cost-versus-accuracy discrepancy. Truth is used only for this post-run
diagnostic, not in initialization or fitting.

These measurements show that allowing all pose variables to move can expose
weak support and imperfect rectangle-surface modeling. Reducing that geometry
cost does not guarantee a more accurate collision center. No threshold was
retuned using the truth diagnostic. Neither a new observer stability proof nor
a deterministic error bound follows from these experiments. Existing
measurement and operating-domain assumptions must still be checked for the
coupled measurement path; native compass heading retains the earlier ideal
heading limitation.

\section{Validation, scope and reproduction}
All 41 Python tests pass, including the existing 32 tests and nine new pose
initialization behaviors. Four focused MATLAB tests pass for current
publication/next-state equivalence, read-only output, analytic maneuver
convergence and pure prediction through dropout. Four actual detector and
MATLAB replays produce 1200 finite estimator publications. The audit records
pose corrections, stage convergence, strict timestamps, exact prefix
invariance, reacquisition, and agreement with cached-detection prototypes.

Another session was editing controller and error-bound files. To keep the
comparison reproducible, this study uses a local copy of committed MATLAB
sources from \path{26cd6f3998d17fd49f1cd9254245805d17f4d3ae}. It reuses the
existing YALMIP/SeDuMi tree by symlink. The initial missing-dependency-path
failure and successful retry are preserved. Concurrent source edits were not
modified or included in this task. Actual inference uses lower scheduling
priority and a 20 percent PyTorch allocation cap on GPU 0; no other process or
CARLA server is changed. Wall times describe offline replay, not a real-time
50 Hz guarantee.

The API is \path{run_pipeline(args,pose_feedback=adapter)}. The adapter supplies
\path{PosePrior(position_ego_m,heading_ego_rad,state_time,last_measurement_time,source)}
from its \texttt{predict} method and receives measured positions or missing
observations through \texttt{update}. Heading-only and pose feedback modes
cannot be selected together. Existing default behavior is preserved.

Source, tests and compact results are tracked. Local replay orchestration,
MATLAB snapshots, caches, raw recordings, model weights and generated figures
remain outside Git. Identified technical exports and command/source hashes
support reproduction without introducing CARLA tooling into the repository.
\IfFileExists{pose_trace.pdf}{\begin{figure}[p]\centering
\includegraphics[width=\textwidth]{pose_trace.pdf}
\caption{Actual perception and estimator errors. The vertical line marks the
fixed 2 s acquisition exclusion; shading marks the camera interruption used
only in the interrupted pose-feedback case.}\end{figure}}{}
\end{document}
